% Kommentare für den Editor (TexWorks/TexMakerX)
% !TeX encoding   = utf8
% !TeX spellcheck = de-DE

% --- LaTeX Vorlage ------------------------------------------------------
% Vorlage für Praktikumsprotokolle in den Modulen Verteilte Systeme, Betriebssysteme
%
% Autor: Aljoscha Pörtner (aljoscha.poertner@fh-bielefeld.de)
% ------------------------------------------------------------------------

% Dokumentenklasse (Koma Script) -----------------------------------------
\documentclass[%
   %draft,     % Entwurfsstadium
   final,      % fertiges Dokument
   paper=a4, paper=portrait, pagesize=auto, % Papier Einstellungen
   fontsize=11pt, % Schriftgröße
   ngerman, % Sprache 
 ]{scrartcl} % Classes: scrartcl, scrreprt, scrbook

% ~~~~~~~~~~~~~~~~~~~~~~~~~~~~~~~~~~~~~~~~~~~~~~~~~~~~~~~~~~~~~~~~~~~~~~~~
% encoding
% ~~~~~~~~~~~~~~~~~~~~~~~~~~~~~~~~~~~~~~~~~~~~~~~~~~~~~~~~~~~~~~~~~~~~~~~~

% Encoding der Dateien (sonst funktionieren Umlaute nicht)
\usepackage[utf8]{inputenc}

% Encoding der Verzeichnisse (für Pfade mit Umlauten und Leerzeichne)
\usepackage[%
   extendedchars, encoding, multidot, space,
   filenameencoding=latin1, % Windows XP, Vista, 7
   % filenameencoding=utf8,   % Linux, OS X
]{grffile}

% ~~~~~~~~~~~~~~~~~~~~~~~~~~~~~~~~~~~~~~~~~~~~~~~~~~~~~~~~~~~~~~~~~~~~~~~~
% Pakete und Stile
% ~~~~~~~~~~~~~~~~~~~~~~~~~~~~~~~~~~~~~~~~~~~~~~~~~~~~~~~~~~~~~~~~~~~~~~~~
% Schriften
% ~~~~~~~~~~~~~~~~~~~~~~~~~~~~~~~~~~~~~~~~~~~~~~~~~~~~~~~~~~~~~~~~~~~~~~~~
% Fonts Fonts Fonts
% ~~~~~~~~~~~~~~~~~~~~~~~~~~~~~~~~~~~~~~~~~~~~~~~~~~~~~~~~~~~~~~~~~~~~~~~~

% immer laden:
\usepackage[T1]{fontenc} % T1 Schrift Encoding
\usepackage{textcomp}	 % Zusätzliche Symbole (Text Companion font extension)

% ~~~~~~~~~~~~~~~~~~~~~~~~~~~~~~~~~~~~~~~~~~~~~~~~~~~~~~~~~~~~~~~~~~~~~~~~
% Symbole
% ~~~~~~~~~~~~~~~~~~~~~~~~~~~~~~~~~~~~~~~~~~~~~~~~~~~~~~~~~~~~~~~~~~~~~~~~

\usepackage{amssymb}
\usepackage{mathcomp}


%% ==== Zusammengesetzte Schriften  (Sans + Serif) =======================

%% - Latin Modern
\usepackage{lmodern}
%% -------------------

%% - Bera Schriften
%\usepackage{bera}
%% -------------------

%% - Times, Helvetica, Courier (Word Standard...)
%\usepackage{mathptmx}
%\usepackage[scaled=.90]{helvet}
%\usepackage{courier}
%% -------------------

%% - Palantino , Helvetica, Courier
%\usepackage{mathpazo}
%\usepackage[scaled=.95]{helvet}
%\usepackage{courier}
%% -------------------

%% - Charter, Bera Sans
%\usepackage{charter}\linespread{1.05}
%\renewcommand{\sfdefault}{fvs}
%\usepackage[charter]{mathdesign}



%%%% =========== Typewriter =============

%\usepackage{courier}                   %% --- Courier
%\renewcommand{\ttdefault}{cmtl}        %% --- CmBright Typewriter Font
%\usepackage[%                          %% --- Luxi Mono (Typewriter)
%   scaled=0.9
%]{luximono}



% Pakete Laden
% ~~~~~~~~~~~~~~~~~~~~~~~~~~~~~~~~~~~~~~~~~~~~~~~~~~~~~~~~~~~~~~~~~~~~~~~~
% These packages must be loaded before all others
% (primarily because they are required by other packages)
% ~~~~~~~~~~~~~~~~~~~~~~~~~~~~~~~~~~~~~~~~~~~~~~~~~~~~~~~~~~~~~~~~~~~~~~~~
\usepackage{calc}
\usepackage{fixltx2e}	% Fix known LaTeX2e bugs

\usepackage[ngerman]{babel} 			% Sprache
\usepackage[dvipsnames, table]{xcolor} 	% Farben

% ~~~~~~~~~~~~~~~~~~~~~~~~~~~~~~~~~~~~~~~~~~~~~~~~~~~~~~~~~~~~~~~~~~~~~~~~
% Bilder, Gleitumgebungen und Platzierung
% ~~~~~~~~~~~~~~~~~~~~~~~~~~~~~~~~~~~~~~~~~~~~~~~~~~~~~~~~~~~~~~~~~~~~~~~~

\usepackage[]{graphicx}					% Graphiken
\usepackage{epstopdf}		% konvertiert eps in pdf

% provides new floats and enables H float modifier option
\usepackage{float}
% Floats immer erst nach der Referenz setzen
\usepackage{flafter}
% Alel Floats werden vor der nächsten section ausgegeben
\usepackage[section]{placeins} 
%

% ~~~~~~~~~~~~~~~~~~~~~~~~~~~~~~~~~~~~~~~~~~~~~~~~~~~~~~~~~~~~~~~~~~~~~~~~
% Beschriftungen (captions)
% ~~~~~~~~~~~~~~~~~~~~~~~~~~~~~~~~~~~~~~~~~~~~~~~~~~~~~~~~~~~~~~~~~~~~~~~~

\usepackage{caption}
\usepackage{subcaption}

% ~~~~~~~~~~~~~~~~~~~~~~~~~~~~~~~~~~~~~~~~~~~~~~~~~~~~~~~~~~~~~~~~~~~~~~~~
% Math
% ~~~~~~~~~~~~~~~~~~~~~~~~~~~~~~~~~~~~~~~~~~~~~~~~~~~~~~~~~~~~~~~~~~~~~~~~

% Base Math Package
\usepackage[fleqn]{amsmath} 
% Warnt bei Benutzung von Befehlen die mit amsmath inkompatibel sind.
\usepackage[all, error]{onlyamsmath}

%Symbols
\usepackage[nointegrals]{wasysym}
% ~~~~~~~~~~~~~~~~~~~~~~~~~~~~~~~~~~~~~~~~~~~~~~~~~~~~~~~~~~~~~~~~~~~~~~~~
% Science
% ~~~~~~~~~~~~~~~~~~~~~~~~~~~~~~~~~~~~~~~~~~~~~~~~~~~~~~~~~~~~~~~~~~~~~~~~

% Einheiten und Zahlenformatierung
\usepackage{siunitx}

% ~~~~~~~~~~~~~~~~~~~~~~~~~~~~~~~~~~~~~~~~~~~~~~~~~~~~~~~~~~~~~~~~~~~~~~~~
% Tables (Tabular)
% ~~~~~~~~~~~~~~~~~~~~~~~~~~~~~~~~~~~~~~~~~~~~~~~~~~~~~~~~~~~~~~~~~~~~~~~~

\usepackage{booktabs}
\usepackage{ltxtable} % Longtable + tabularx

% ~~~~~~~~~~~~~~~~~~~~~~~~~~~~~~~~~~~~~~~~~~~~~~~~~~~~~~~~~~~~~~~~~~~~~~~~
% text related packages
% ~~~~~~~~~~~~~~~~~~~~~~~~~~~~~~~~~~~~~~~~~~~~~~~~~~~~~~~~~~~~~~~~~~~~~~~~

\usepackage{url}            % Befehl \url{...}
\usepackage{enumitem}		% Kompakte Listen

% Neue Befehle: \Centering, \RaggedLeft, and \RaggedRight, ... 
\usepackage{ragged2e}


% ~~~~~~~~~~~~~~~~~~~~~~~~~~~~~~~~~~~~~~~~~~~~~~~~~~~~~~~~~~~~~~~~~~~~~~~~
% Citations
% ~~~~~~~~~~~~~~~~~~~~~~~~~~~~~~~~~~~~~~~~~~~~~~~~~~~~~~~~~~~~~~~~~~~~~~~~

%\usepackage[
%	style=alphabetic, % Loads the bibliography and the citation style 
%	natbib=true, % define natbib compatible cite commands
%]{biblatex}	
% Other options:
%	style=numeric, % 
%	style=numeric-comp,    % [1–3, 7, 8]
%	style=numeric-verb,    % [2]; [5]; [6]


% ~~~~~~~~~~~~~~~~~~~~~~~~~~~~~~~~~~~~~~~~~~~~~~~~~~~~~~~~~~~~~~~~~~~~~~~~
% layout packages
% ~~~~~~~~~~~~~~~~~~~~~~~~~~~~~~~~~~~~~~~~~~~~~~~~~~~~~~~~~~~~~~~~~~~~~~~~
%
% Befehle für 1,5 und 2 zeilig: 
% \singlespacing, \onehalfspacing und \doublespacing
\usepackage{setspace}

% ~~~~~~~~~~~~~~~~~~~~~~~~~~~~~~~~~~~~~~~~~~~~~~~~~~~~~~~~~~~~~~~~~~~~~~~~
% Kopf und Fusszeile
% ~~~~~~~~~~~~~~~~~~~~~~~~~~~~~~~~~~~~~~~~~~~~~~~~~~~~~~~~~~~~~~~~~~~~~~~~

% Kopf und Fusszeile mit scrpage2 einstellen
\usepackage[automark, komastyle, nouppercase]{scrpage2}

% ~~~~~~~~~~~~~~~~~~~~~~~~~~~~~~~~~~~~~~~~~~~~~~~~~~~~~~~~~~~~~~~~~~~~~~~~
% pdf packages
% ~~~~~~~~~~~~~~~~~~~~~~~~~~~~~~~~~~~~~~~~~~~~~~~~~~~~~~~~~~~~~~~~~~~~~~~~

% Include pages from external PDF documents in LaTeX documents
\usepackage{pdfpages} 

% Optischer Randausgleich mit pdfTeX
\usepackage{microtype}

\usepackage[unicode]{hyperref}

\usepackage{listings}


\usepackage[style=ieee,backend=biber]{biblatex}
% Einstellungen und Layoutstile 
% ~~~~~~~~~~~~~~~~~~~~~~~~~~~~~~~~~~~~~~~~~~~~~~~~~~~~~~~~~~~~~~~~~~~~~~~~
% Colors
% ~~~~~~~~~~~~~~~~~~~~~~~~~~~~~~~~~~~~~~~~~~~~~~~~~~~~~~~~~~~~~~~~~~~~~~~~
\definecolor{sectioncolor}{RGB}{0, 0, 0}     % black

% ~~~~~~~~~~~~~~~~~~~~~~~~~~~~~~~~~~~~~~~~~~~~~~~~~~~~~~~~~~~~~~~~~~~~~~~~
% text related 
% ~~~~~~~~~~~~~~~~~~~~~~~~~~~~~~~~~~~~~~~~~~~~~~~~~~~~~~~~~~~~~~~~~~~~~~~~

%% style of URL
\urlstyle{tt}


% Keine hochgestellten Ziffern in der Fussnote (KOMA-Script-spezifisch):
\deffootnote{1.5em}{1em}{\makebox[1.5em][l]{\thefootnotemark}}

% Limit space of footnotes to 10 lines
\setlength{\dimen\footins}{10\baselineskip}

% prevent continuation of footnotes 
% at facing page
\interfootnotelinepenalty=10000 

% ~~~~~~~~~~~~~~~~~~~~~~~~~~~~~~~~~~~~~~~~~~~~~~~~~~~~~~~~~~~~~~~~~~~~~~~~
% Science
% ~~~~~~~~~~~~~~~~~~~~~~~~~~~~~~~~~~~~~~~~~~~~~~~~~~~~~~~~~~~~~~~~~~~~~~~~

\sisetup{%
	mode = math, detect-family, detect-weight,	
	exponent-product = \cdot,
	number-unit-separator=\text{\,},
	output-decimal-marker={,},
}

% ~~~~~~~~~~~~~~~~~~~~~~~~~~~~~~~~~~~~~~~~~~~~~~~~~~~~~~~~~~~~~~~~~~~~~~~~
% Citations / Style of Bibliography
% ~~~~~~~~~~~~~~~~~~~~~~~~~~~~~~~~~~~~~~~~~~~~~~~~~~~~~~~~~~~~~~~~~~~~~~~~

% Kommentar entfernene wenn biblatex geladen wird
% \IfPackageLoaded{biblatex}{%
	\ExecuteBibliographyOptions{%
%--- Backend --- --- ---
	backend=bibtex,  % (bibtex, bibtex8, biber)
	bibwarn=true, %
	bibencoding=ascii, % (ascii, inputenc, <encoding>)
%--- Sorting --- --- ---
	sorting=nty, % Sort by name, title, year.
	% other options: 
	% nty        Sort by name, title, year.
	% nyt        Sort by name, year, title.
	% nyvt       Sort by name, year, volume, title.
	% anyt       Sort by alphabetic label, name, year, title.
	% anyvt      Sort by alphabetic label, name, year, volume, title.
	% ynt        Sort by year, name, title.
	% ydnt       Sort by year (descending), name, title.
	% none       Do not sort at all. All entries are processed in citation order.
	% debug      Sort by entry key. This is intended for debugging only.
	%
	sortcase=true,
	sortlos=los, % (bib, los) The sorting order of the list of shorthands
	sortcites=false, % do/do not sort citations according to bib	
%--- Dates --- --- ---
	date=comp,  % (short, long, terse, comp, iso8601)
%	origdate=
%	eventdate=
%	urldate=
%	alldates=
	datezeros=true, %
	dateabbrev=true, %
%--- General Options --- --- ---
	maxnames=1,
	minnames=1,
%	maxbibnames=99,
%	maxcitenames=1,
%	autocite= % (plain, inline, footnote, superscript) 
	autopunct=true,
	language=auto,
	babel=none, % (none, hyphen, other, other*)
	block=none, % (none, space, par, nbpar, ragged)
	notetype=foot+end, % (foot+end, footonly, endonly)
	hyperref=true, % (true, false, auto)
	backref=true,
	backrefstyle=three, % (none, three, two, two+, three+, all+)
	backrefsetstyle=setonly, %
	indexing=false, % 
	% options:
	% true       Enable indexing globally.
	% false      Disable indexing globally.
	% cite       Enable indexing in citations only.
	% bib        Enable indexing in the bibliography only.
	refsection=none, % (part, chapter, section, subsection)
	refsegment=none, % (none, part, chapter, section, subsection)
	abbreviate=true, % (true, false)
	defernumbers=false, % 
	punctfont=false, % 
	arxiv=abs, % (ps, pdf, format)	
%--- Style Options --- --- ---	
% The following options are provided by the standard styles
	isbn=false,%
	url=false,%
	doi=false,%
	eprint=false,%	
	}%	
	
	% change alpha label to be without +	
	\renewcommand*{\labelalphaothers}{}
	
	% change 'In: <magazine>" to "<magazine>"
	\renewcommand*{\intitlepunct}{}
	\DefineBibliographyStrings{german}{in={}}
	
	% make names capitalized \textsc{}
	\renewcommand{\mkbibnamefirst}{\textsc}
	\renewcommand{\mkbibnamelast}{\textsc}
	
	% make volume and number look like 
	% 'Bd. 33(14): '
	\renewbibmacro*{volume+number+eid}{%
	  \setunit{\addcomma\space}%
	  \bibstring{volume}% 
	  \setunit{\addspace}%
	  \printfield{volume}%
	  \iffieldundef{number}{}{% 
	    \printtext[parens]{%
	      \printfield{number}%
	    }%
	  }%
	  \setunit{\addcomma\space}%
	  \printfield{eid}
	  %\setunit{\addcolon\space}%
	  }	

	% <authors>: <title>
	\renewcommand*{\labelnamepunct}{\addcolon\space}
	% make ': ' before pages
	\renewcommand*{\bibpagespunct}{\addcolon\space}
	% names delimiter ';' instead of ','
	%\renewcommand*{\multinamedelim}{\addsemicolon\space}

	% move date before issue
	\renewbibmacro*{journal+issuetitle}{%
	  \usebibmacro{journal}%
	  \setunit*{\addspace}%
	  \iffieldundef{series}
	    {}
	    {\newunit
	     \printfield{series}%
	     \setunit{\addspace}}%
	  %
	  \usebibmacro{issue+date}%
	  \setunit{\addcolon\space}%
	  \usebibmacro{issue}%
	  \setunit{\addspace}%
	  \usebibmacro{volume+number+eid}%
	  \newunit}

	% print all names, even if maxnames = 1
	\DeclareCiteCommand{\citeauthors}
	  {
	   \defcounter{maxnames}{1000}
	   \boolfalse{citetracker}%
	   \boolfalse{pagetracker}%
	   \usebibmacro{prenote}}
	  {\ifciteindex
	     {\indexnames{labelname}}
	     {}%
	   \printnames{labelname}}
	  {\multicitedelim}
	  {\usebibmacro{postnote}}

}%

% ~~~~~~~~~~~~~~~~~~~~~~~~~~~~~~~~~~~~~~~~~~~~~~~~~~~~~~~~~~~~~~~~~~~~~~~~
% figures, placement, floats and captions
% ~~~~~~~~~~~~~~~~~~~~~~~~~~~~~~~~~~~~~~~~~~~~~~~~~~~~~~~~~~~~~~~~~~~~~~~~

% Make float placement easier
\renewcommand{\floatpagefraction}{.75} % vorher: .5
\renewcommand{\textfraction}{.1}       % vorher: .2
\renewcommand{\topfraction}{.8}        % vorher: .7
\renewcommand{\bottomfraction}{.5}     % vorher: .3
\setcounter{topnumber}{3}        % vorher: 2
\setcounter{bottomnumber}{2}     % vorher: 1
\setcounter{totalnumber}{5}      % vorher: 3

%% ~~~ Captions ~~~~~~~~~~~~~~~~~~~~~~~~~~~~~~~~~~~~~~~~~~~~~~~~~~~~~~~~~~
% Style of captions
\DeclareCaptionStyle{captionStyleTemplateDefault}
[ % single line captions
   justification = centering
]
{ % multiline captions
% -- Formatting
   format      = plain,  % plain, hang
   indention   = 0em,    % indention of text 
   labelformat = default,% default, empty, simple, brace, parens
   labelsep    = colon,  % none, colon, period, space, quad, newline, endash
   textformat  = simple, % simple, period
% -- Justification
   justification = justified, %RaggedRight, justified, centering
   singlelinecheck = true, % false (true=ignore justification setting in single line)
% -- Fonts
   labelfont   = {small,bf},
   textfont    = {small,rm},
% valid values:
% scriptsize, footnotesize, small, normalsize, large, Large
% normalfont, ip, it, sl, sc, md, bf, rm, sf, tt
% singlespacing, onehalfspacing, doublespacing
% normalcolor, color=<...>
%
% -- Margins and further paragraph options
   margin = 10pt, %.1\textwidth,
   % width=.8\linewidth,
% -- Skips
   skip     = 10pt, % vertical space between the caption and the figure
   position = auto, % top, auto, bottom
% -- Lists
   % list=no, % suppress any entry to list of figure 
   listformat = subsimple, % empty, simple, parens, subsimple, subparens
% -- Names & Numbering
   % figurename = Abb. %
   % tablename  = Tab. %
   % listfigurename=
   % listtablename=
   % figurewithin=chapter
   % tablewithin=chapter
%-- hyperref related options
	hypcap=true, % (true, false) 
	% true=all hyperlink anchors are placed at the 
	% beginning of the (floating) environment
	%
	hypcapspace=0.5\baselineskip
}

% apply caption style
\captionsetup{
	style = captionStyleTemplateDefault % base
}

% Predefinded skip setup for different floats
\captionsetup[table]{position=top}
\captionsetup[figure]{position=bottom}


% options for subcaptions
\captionsetup[sub]{ %
	style = captionStyleTemplateDefault, % base
	skip=6pt,
	margin=5pt,
	labelformat = parens,% default, empty, simple, brace
	labelsep    = space,
	list=false,
	hypcap=false
}

% ~~~~~~~~~~~~~~~~~~~~~~~~~~~~~~~~~~~~~~~~~~~~~~~~~~~~~~~~~~~~~~~~~~~~~~~~
% layout 
% ~~~~~~~~~~~~~~~~~~~~~~~~~~~~~~~~~~~~~~~~~~~~~~~~~~~~~~~~~~~~~~~~~~~~~~~~


%% Paragraph Separation =================================
\KOMAoptions{%
   parskip=absolute, % do not change indentation according to fontsize
   parskip=false     % indentation of 1em
   % parskip=half    % parksip of 1/2 line 
}%

%% line spacing =========================================
%\onehalfspacing	% 1,5-facher Abstand
%\doublespacing		% 2-facher Abstand

%% page layout ==========================================

\raggedbottom     % Variable Seitenhoehen zulassen

% Koma Script text area layout
\KOMAoptions{%
   DIV=11,% (Size of Text Body, higher values = greater textbody)
   BCOR=5mm% (Bindekorrektur)
}%

%%% === Page Layout  Options ===
\KOMAoptions{% (most options are for package typearea)
   % twoside=true, % two side layout (alternating margins, standard in books)
   twoside=false, % single side layout 
   %
   headlines=2.1,%
}%

%\KOMAoptions{%
%      headings=noappendixprefix % chapter in appendix as in body text
%      ,headings=nochapterprefix  % no prefix at chapters
%      % ,headings=appendixprefix   % inverse of 'noappendixprefix'
%      % ,headings=chapterprefix    % inverse of 'nochapterprefix'
%      % ,headings=openany   % Chapters start at any side
%      % ,headings=openleft  % Chapters start at left side
%      ,headings=openright % Chapters start at right side      
%}%


% reloading of typearea, necessary if setting of spacing changed
\typearea[current]{last}

% ~~~~~~~~~~~~~~~~~~~~~~~~~~~~~~~~~~~~~~~~~~~~~~~~~~~~~~~~~~~~~~~~~~~~~~~~
% Titlepage
% ~~~~~~~~~~~~~~~~~~~~~~~~~~~~~~~~~~~~~~~~~~~~~~~~~~~~~~~~~~~~~~~~~~~~~~~~
\KOMAoptions{%
   % titlepage=true %
   titlepage=false %
}%

% ~~~~~~~~~~~~~~~~~~~~~~~~~~~~~~~~~~~~~~~~~~~~~~~~~~~~~~~~~~~~~~~~~~~~~~~~
% head and foot lines
% ~~~~~~~~~~~~~~~~~~~~~~~~~~~~~~~~~~~~~~~~~~~~~~~~~~~~~~~~~~~~~~~~~~~~~~~~

% \pagestyle{scrheadings} % Seite mit Headern
\pagestyle{scrplain} % Seiten ohne Header

% loescht voreingestellte Stile
\clearscrheadings
\clearscrplain
%
% Was steht wo...
% Bei headings:
%   % Oben aussen: Kapitel und Section
%   % Unten aussen: Seitenzahl
%   \ohead{\pagemark}
%   \ihead{\headmark}
%   \ofoot[\pagemark]{} % Außen unten: Seitenzahlen bei plain
% Bei Plain:
\cfoot[\pagemark]{\pagemark} % Mitte unten: Seitenzahlen bei plain


% Angezeigte Abschnitte im Header
% \automark[section]{chapter} %[rechts]{links}
\automark[subsection]{section} %[rechts]{links}

% ~~~~~~~~~~~~~~~~~~~~~~~~~~~~~~~~~~~~~~~~~~~~~~~~~~~~~~~~~~~~~~~~~~~~~~~~
% headings / page opening
% ~~~~~~~~~~~~~~~~~~~~~~~~~~~~~~~~~~~~~~~~~~~~~~~~~~~~~~~~~~~~~~~~~~~~~~~~
\setcounter{secnumdepth}{2}

\KOMAoptions{%
%%%% headings
   % headings=small  % Small Font Size, thin spacing above and below
   % headings=normal % Medium Font Size, medium spacing above and below
   headings=big % Big Font Size, large spacing above and below
}%

% Titelzeile linksbuendig, haengend
\renewcommand*{\raggedsection}{\raggedright} 

% ~~~~~~~~~~~~~~~~~~~~~~~~~~~~~~~~~~~~~~~~~~~~~~~~~~~~~~~~~~~~~~~~~~~~~~~~
% fonts of headings
% ~~~~~~~~~~~~~~~~~~~~~~~~~~~~~~~~~~~~~~~~~~~~~~~~~~~~~~~~~~~~~~~~~~~~~~~~
\setkomafont{sectioning}{\normalfont\sffamily} % \rmfamily
\setkomafont{descriptionlabel}{\itshape}
\setkomafont{pageheadfoot}{\normalfont\normalcolor\small\sffamily}
\setkomafont{pagenumber}{\normalfont\sffamily}

%%% --- Titlepage ---
%\setkomafont{subject}{}
%\setkomafont{subtitle}{}
%\setkomafont{title}{}

% ~~~~~~~~~~~~~~~~~~~~~~~~~~~~~~~~~~~~~~~~~~~~~~~~~~~~~~~~~~~~~~~~~~~~~~~~
% settings and layout of TOC, LOF, 
% ~~~~~~~~~~~~~~~~~~~~~~~~~~~~~~~~~~~~~~~~~~~~~~~~~~~~~~~~~~~~~~~~~~~~~~~~
\setcounter{tocdepth}{3} % Depth of TOC Display

% ~~~~~~~~~~~~~~~~~~~~~~~~~~~~~~~~~~~~~~~~~~~~~~~~~~~~~~~~~~~~~~~~~~~~~~~~
% Tabellen
% ~~~~~~~~~~~~~~~~~~~~~~~~~~~~~~~~~~~~~~~~~~~~~~~~~~~~~~~~~~~~~~~~~~~~~~~~

%%% -| Neue Spaltendefinitionen 'columntypes' |--
%
% Belegte Spaltentypen:
% l - links
% c - zentriert
% r - rechts
% p,m,b  - oben, mittig, unten
% X - tabularx Auto-Spalte

% um Tabellenspalten mit Flattersatz zu setzen, muss \\ vor
% (z.B.) \raggedright geschuetzt werden:
\newcommand{\PreserveBackslash}[1]{\let\temp=\\#1\let\\=\temp}

% Spalten mit Flattersatz und definierte Breite:
% m{} -> mittig
% p{} -> oben
% b{} -> unten
%
% Linksbuendig:
\newcolumntype{v}[1]{>{\PreserveBackslash\RaggedRight\hspace{0pt}}p{#1}}
\newcolumntype{M}[1]{>{\PreserveBackslash\RaggedRight\hspace{0pt}}m{#1}}
% % Rechtsbuendig :
% \newcolumntype{R}[1]{>{\PreserveBackslash\RaggedLeft\hspace{0pt}}m{#1}}
% \newcolumntype{S}[1]{>{\PreserveBackslash\RaggedLeft\hspace{0pt}}p{#1}}
% % Zentriert :
% \newcolumntype{Z}[1]{>{\PreserveBackslash\Centering\hspace{0pt}}m{#1}}
% \newcolumntype{A}[1]{>{\PreserveBackslash\Centering\hspace{0pt}}p{#1}}

\newcolumntype{Y}{>{\PreserveBackslash\RaggedLeft\hspace{0pt}}X}

%-- Einstellungen für Tabellen ----------
\providecommand\tablestyle{%
  \renewcommand{\arraystretch}{1.4} % Groessere Abstaende zwischen Zeilen
  \normalfont\normalsize            %
  \sffamily\small           % Serifenlose und kleine Schrift
  \centering%                       % Tabelle zentrieren
}

%--Einstellungen für Tabellen ----------

\colorlet{tablesubheadcolor}{gray!40}
\colorlet{tableheadcolor}{gray!25}
\colorlet{tableblackheadcolor}{black!60}
\colorlet{tablerowcolor}{gray!15.0}

\usepackage{color}

\definecolor{pblue}{rgb}{0.13,0.13,1}
\definecolor{pgreen}{rgb}{0,0.5,0}
\definecolor{pred}{rgb}{0.9,0,0}
\definecolor{pgrey}{rgb}{0.46,0.45,0.48}

\usepackage{listings}
\lstset{language=Java,
	showspaces=false,
	showtabs=false,
	breaklines=true,
	frame=L,
	numbers=left,  
	captionpos=b,
	backgroundcolor=\color{white},   % choose the background color; you must add \usepackage{color} or \usepackage{xcolor}
	showstringspaces=false,
	breakatwhitespace=true,
	commentstyle=\color{pgreen},
	keywordstyle=\color{pblue},
	stringstyle=\color{pred},
	basicstyle=\footnotesize,
	moredelim=[il][\textcolor{pgrey}]{$$},
	moredelim=[is][\textcolor{pgrey}]{\%\%}{\%\%}
}
% ~~~~~~~~~~~~~~~~~~~~~~~~~~~~~~~~~~~~~~~~~~~~~~~~~~~~~~~~~~~~~~~~~~~~~~~~
% pdf packages
% ~~~~~~~~~~~~~~~~~~~~~~~~~~~~~~~~~~~~~~~~~~~~~~~~~~~~~~~~~~~~~~~~~~~~~~~~

% ~~~~~~~~~~~~~~~~~~~~~~~~~~~~~~~~~~~~~~~~~~~~~~~~~~~~~~~~~~~~~~~~~~~~~~~~
% fix remaining problems
% ~~~~~~~~~~~~~~~~~~~~~~~~~~~~~~~~~~~~~~~~~~~~~~~~~~~~~~~~~~~~~~~~~~~~~~~~



%\addbibresource{appendix/IEEEabrv.bib}
\addbibresource{appendix/literature.bib}
% ~~~~~~~~~~~~~~~~~~~~~~~~~~~~~~~~~~~~~~~~~~~~~~~~~~~~~~~~~~~~~~~~~~~~~~~~
% Eigene Befehle
% ~~~~~~~~~~~~~~~~~~~~~~~~~~~~~~~~~~~~~~~~~~~~~~~~~~~~~~~~~~~~~~~~~~~~~~~~
% -- new commands --
\providecommand{\abs}[1]{\lvert#1\rvert}
\providecommand{\Abs}[1]{\left\lvert#1\right\rvert}
\providecommand{\norm}[1]{\left\Vert#1\right\Vert}
\providecommand{\Trace}[1]{\ensuremath{\Tr\{\,#1\,\}}} % Trace /Spur
%

\renewcommand{\d}{\partial\mspace{2mu}} % partial diff
\newcommand{\td}{\,\mathrm{d}}	% total diff

\newcommand{\Ham}{\mathcal{H}}    % Hamilton
\newcommand{\Prob}{\mathscr{P}}    % Hamilton
\newcommand{\unity}{\mathds{1}}   % Real

\renewcommand{\i}{\mathrm{i}}   % imagin�re Einheit



% -- New Operators --
\DeclareMathOperator{\rot}{rot}
\DeclareMathOperator{\grad}{grad}
\DeclareMathOperator{\Tr}{Tr}
\DeclareMathOperator{\const}{const}
\DeclareMathOperator{\e}{e} 			% exponatial Function



% ~~~~~~~~~~~~~~~~~~~~~~~~~~~~~~~~~~~~~~~~~~~~~~~~~~~~~~~~~~~~~~~~~~~~~~~~
% Eigene Befehle
% ~~~~~~~~~~~~~~~~~~~~~~~~~~~~~~~~~~~~~~~~~~~~~~~~~~~~~~~~~~~~~~~~~~~~~~~~
% Silbentrennung hinzufügen als 
% Sil-ben-tren-nung 
\hyphenation{}

\listfiles % schreibt alle verwendeten Dateien in die log Datei

%% Dokument Beginn %%%%%%%%%%%%%%%%%%%%%%%%%%%%%%%%%%%%%%%%%%%%%%%%%%%%%%%%
\begin{document}

\input{content/titel}
\input{content/Gelaufgaben}
\tableofcontents

% Testdokumente (auskommentieren!)
\newcommand{\env}[1]{\texttt{#1}}
\newcommand{\command}[1]{\texttt{#1}}
\newcommand{\package}[1]{\texttt{\itshape#1}}
\newcommand{\engl}[1]{(engl: \textit{#1})\xspace}
\newpage

\section*{\LARGE Hinweis zur Erstellung von Praktikumsprotokollen} 
Die Erstellung von Praktikumsprotokollen dient der Vorbereitung sowohl auf die akademische als auch auf die praktische Arbeit. Daher sind die gängigen Regeln des wissenschaftlichen Arbeits einzuhalten. Im Protokoll soll jede Phase des Versuchs bzw. des Praktikums festgehalten und entsprechend dokumentiert werden. Im Folgenden werden die wichtigen Teilaspekte dargestellt und besprochen.

\subsection*{Deckblatt}

Verwenden Sie am Besten, dass bereits vorbereitete Deckblatt. Sollten Sie ein eigenens Deckblatt bevorzugen, sind mindestens die folgenden Einträge zu übernehmen:

\begin{itemize}
	\item Praktikumstitel und Praktikumsnummer
	\item Datum
	\item Versuchsteilnehmer mit Matrikelnummer
	\item Protokollführer (Erste Stelle der Versuchsteilnehmer)
\end{itemize}

\subsection*{Frage- bzw. Aufgabenstellung}
Für jede Teilaufgabe ist die bearbeitete Frage- bzw. Aufgabenstellung unbedingt mit zu übertragen. (Nicht nur die Nummer !)

\subsubsection{Lösungsweg}
Beschreiben Sie die Durchführung der Aufgabenstellung. Hieraus folgt, dass Sie nicht nur die Ergebnisse festhalten, sondern auch den Weg dahin und zwar so kleinschrittig wie nötig und \underline{sinnvoll}. Aufgetretene Effekte bzw. Probleme sind ebenso festzuhalten. Ziel ist es den Versuch für Dritte nachvollziehbar zu gestalten.

\subsubsection{Ergebnisse}
Beschreiben Sie in jedem Fall die erzielten Ergebnisse. Selbst wenn ein Versuch nicht das gewünschte Ergebnis erziehlt hat, sollten Sie diese trotzdem dokumentieren. (\textbf{Für nicht dokumentierte Ergebnisse können keine Punkte vergeben werden})

\subsubsection{Diskussion}
\begin{itemize}
	\item In jedem Fall sollte innerhalb der Diskussion auf die Fragestellung eingegangen werden, jedoch sind auch weitere Aspekte des Versuchs zu diskutieren.
	\newpage
	\item Probleme und besondere Vorkomnisse: Probleme können auftreteten, sollten jedoch in jedem Fall im Protokoll festgehalten werden. Neben der Meldung bei dem entsprechenden Betreuer sind folgende Punkte zu dokumentieren:
	\begin{itemize}
		\item was aufgetreten ist und unter welchen Bedingungen
		\item was Sie versucht haben und mit welchen Ergebnissen
		\item die möglichen Ursachen des Problems
	\end{itemize}
\end{itemize}

\subsubsection{Rückmeldung}
Wir würden uns freuen, wenn Sie das Protokoll auch durchaus nutzen, um Rückmeldung zu den Praktikas zu geben. Gab es Probleme bei der Hardware, ist die Organisation verbesserungswürdig oder gab es unklare Fragestellungen. \\

\textbf{Hinweis}: Sollten wir den Eindruck erlangen, dass Sie das Protokoll nicht eigenständig (in Ihrer Gruppe) angefertigt haben oder dass das Protokoll auf Basis von den Ergebnissen anderer Gruppen angefertigt wurde, wird das entsprechende Praktikum mit 5.0 bewertet und die Abgabe als Täuschungsversuch gesehen.
\\

\textbf{Hinweis:} Neben der Abgabe des Protokolls ist die Teilnahme an Praktikas weiterhin Pflicht und eine notwendige Teilleistung für das Bestehen des Praktikums. In diesen werden Sie Ihre Ergebnisse vorstellen und wir werden diese besprechen.

\newpage
\section{Checkliste für das Praktikumsprotokoll}

\begin{itemize}
	\item[\Square] \textbf{Struktur und Aufbau des Protokolls eingehalten}
	\begin{itemize}
		\item Inhaltsverzeichnis ?
		\item Tabelle mit gelösten und nicht gelösten Aufgaben ?
		\item Für jede Aufgabe: Vorbereitung, Lösungsweg, Ergebnis ?
	\end{itemize}
	\item[\Square] \textbf{Das Protokoll enthält die gesamte, klar formulierte Lösung}
	\begin{itemize}
		\item Vorbedingungen ?
		\item Sourcecode ?
		\item Screenshots ?
		\item Formeln ?
		\item Rechenweg?
		\item Evaluation ?
		\item Erklärungen ?
	\end{itemize}
	\item[\Square] \textbf{Sourcecode als Codesnippets mit erklärendem Text}
	\begin{itemize}
		\item Keine eigene Seite für Sourcecode (Ausnahme: Anhang) ?
		\item Wichtige Codezeilen beschrieben ?
		\item Listings benutzt ?
	\end{itemize}
	\item[\Square] \textbf{Nachweis der Quellen}
	\begin{itemize}
		\item Literaturquellen angegeben und richtig zitiert ?
		\item Wenn Beispielsourcecode verwendet (bspw. aus dem Internet), ist dieser klar erkennbar und referenziert ?
	\end{itemize}
\end{itemize}

\newpage

\paragraph{\LARGE Struktur eines Beispielprotokolls}

\section{Aufgabe 1}
	\subsection{Vorbereitung}
	Was für Vorbereitungen waren nötig, um die Aufgabe zu lösen. Hier können beispielsweise Konfigurationen von Tools oder VM's beschrieben werden.
	\subsection{Durchführung}
	Wie ist die konkrete Durchführung vonstattengegangen. Wichtig hierbei ist die detaillierte Beschreibung des ``Versuches'' eventuell auch unter Zuhilfenahme von Diagrammen oder Listings.
	\subsection{Fazit}
	Hat der Versuch geklappt? Welche Probleme traten auf? Welche Ergebnisse wurden erzielt. Wie lassen sich die Ergebnisse interpretieren.
	
\begin{quote}
\textit{Hinweis: Wichtig ist, dass bei der Dokumentation der Ergebnisse sämtliche Quellen, die bei der Erfüllung der Aufgabe hilfreich waren, korrekt zitiert werden. (Stackoverflow ist keine (alleinige) Quelle und sollte bei Bedarf durch weitere Literatur untermauert werden - Warum funktioniert die Lösung aus einem Forum etc.).}
\end{quote}
\section{Aufgabe 2}
	\subsection{Vorbereitung}
	...
	\subsection{Durchführung}
	...
	\subsection{Fazit}
	...
\newpage
\paragraph{{\LARGE Demonstration eines Dokumentes (nach Matthias Pospiech)}}

% ============================================================
\section{Textauszeichnungen}

Die Standard Text Schalter erzeugen
\textbf{fett},
\textit{kursiv\footnote{kursiv und schräggestellt sind unterschiedliche Schriften und daher nicht das gleiche !} (italic)},
\textsl{schräggestellt (slanted)},
\textsf{serifenlos (grotesk)},
\textsc{Kapitälchen} und
\texttt{Schreibmaschinenschrift}.
Sowie beliebige Kombinationen derselben:
\textit{\textbf{fett kursiv}},
\textsl{\textbf{fett schräg}},
\textsf{\textbf{fett serifenlos}},
\textsc{\textbf{Fette Kapitälchen}}
und
\textsl{\textsf{serifenlos schräg}}. Je nach Schrift sind jedoch nicht alle Kombinationen möglich. In dem Fall bekommt man die Fehlermeldung `Some font shapes were not available, defaults substituted.'

% ============================================================

\section{Text und Zitate}
% ------------------------------------------------------------
\subsection{Linksbündig}
%
Mit dem Paket \package{ragged2e} und der Umgebung \env{FlushLeft} bzw. dem Befehl \command{Raggedright} wird Flattersatz (Gegenteil von Blocksatz) gegenüber dem \LaTeX{} Standard noch verbessert.

\begin{FlushLeft}
Aliquam ultrices libero hendrerit diam. Vestibulum ultrices sapien sit amet elit. Quisque tempor nisl eu sem. Nam lorem lectus, viverra nec, rutrum quis, lobortis nec, magna. Praesent hendrerit tortor vitae elit. Vivamus sed leo at mi elementum semper. Lorem ipsum dolor sit amet, consectetuer adipiscing elit. Aliquam eu nisi. Nam eget dui a tortor congue imperdiet. Etiam mattis. Nam tristique. Sed malesuada neque ut leo. Aenean est. In id augue.
\end{FlushLeft}
% ------------------------------------------------------------
\subsection{Rechtsbündig}
%
Mit dem Paket \package{ragged2e} und der Umgebung \env{FlushRight} bzw. dem Befehl \command{RaggedLeft} wird Flattersatz (Gegenteil von Blocksatz) gegenüber dem \LaTeX{} Standard noch verbessert.

\begin{FlushRight}
Aliquam ultrices libero hendrerit diam. Vestibulum ultrices sapien sit amet elit. Quisque tempor nisl eu sem. Nam lorem lectus, viverra nec, rutrum quis, lobortis nec, magna. Praesent hendrerit tortor vitae elit. Vivamus sed leo at mi elementum semper. Lorem ipsum dolor sit amet, consectetuer adipiscing elit. Aliquam eu nisi. Nam eget dui a tortor congue imperdiet. Etiam mattis. Nam tristique. Sed malesuada neque ut leo. Aenean est. In id augue.
\end{FlushRight}
%
% ------------------------------------------------------------
\subsection{Zentierter Text}
%
Zentierter Text, erstellt mit \env{center}\footnote{nicht verwenden innerhalb \env{table} oder \env{figure}!}.
\begin{center}
Aliquam ultrices libero hendrerit diam. Vestibulum ultrices sapien sit amet elit. Quisque tempor nisl eu sem. Nam lorem lectus, viverra nec, rutrum quis, lobortis nec, magna. Praesent hendrerit tortor vitae elit. Vivamus sed leo at mi elementum semper. Lorem ipsum dolor sit amet, consectetuer adipiscing elit. Aliquam eu nisi. Nam eget dui a tortor congue imperdiet. Etiam mattis. Nam tristique. Sed malesuada neque ut leo. Aenean est. In id augue.
\end{center}

% ------------------------------------------------------------
\subsection{Zitat mit quote}
Ein Text in einer \env{quote} Umgebung\footnote{Für Zitate eignet sich allerdings eher die Umgebung \env{blockquote} vom Paket \package{csquotes}}. Der gesammte Text innerhalb der Umgebung wird von beiden Seiten eingerückt.

\begin{quote}
Aliquam ultrices libero hendrerit diam. Vestibulum ultrices sapien sit amet elit. Quisque tempor nisl eu sem. Nam lorem lectus, viverra nec, rutrum quis, lobortis nec, magna. Praesent hendrerit tortor vitae elit. Vivamus sed leo at mi elementum semper. Lorem ipsum dolor sit amet, consectetuer adipiscing elit. Aliquam eu nisi. Nam eget dui a tortor congue imperdiet. Etiam mattis. Nam tristique. Sed malesuada neque ut leo. Aenean est. In id augue.
\end{quote}



% ============================================================
\section{Listen}

% ------------------------------------------------------------
\subsection{itemize}
Dies ist die Standard Aufzählungsliste von \LaTeXe. Sie hat einen Abstand zwischen den Einträgen um bei Längeren Zeilen das Lesen zu erleichtern.

\begin{itemize}
   \item Lorem ipsum dolor sit amet, consectetur adipisicing elit, sed do eiusmod tempor incididunt ut labore et dolore magna aliqua.
%
   \item Lorem ipsum dolor sit amet, consectetur adipisicing elit, sed do eiusmod tempor incididunt ut labore et dolore magna aliqua.
%
   \item Lorem ipsum dolor sit amet, consectetur adipisicing elit, sed do eiusmod tempor incididunt ut labore et dolore magna aliqua.
\end{itemize}

mit mehr als einer Ebene

\begin{itemize}
   \item Lorem ipsum dolor sit amet, consectetur adipisicing elit, sed do eiusmod tempor incididunt ut labore et dolore magna aliqua.
%
   \begin{itemize}
      \item Lorem ipsum dolor sit amet, consectetur adipisicing elit, sed do eiusmod tempor incididunt ut labore et dolore magna aliqua.
      %
      \begin{itemize}
         \item Lorem ipsum dolor sit amet, consectetur adipisicing elit, sed do eiusmod tempor incididunt ut labore et dolore magna aliqua.
         %
         \item Lorem ipsum dolor sit amet, consectetur adipisicing elit, sed do eiusmod tempor incididunt ut labore et dolore magna aliqua.
         %
         \item Lorem ipsum dolor sit amet, consectetur adipisicing elit, sed do eiusmod tempor incididunt ut labore et dolore magna aliqua.
      \end{itemize}
      %
      \item Lorem ipsum dolor sit amet, consectetur adipisicing elit, sed do eiusmod tempor incididunt ut labore et dolore magna aliqua.
      %
      \item Lorem ipsum dolor sit amet, consectetur adipisicing elit, sed do eiusmod tempor incididunt ut labore et dolore magna aliqua.
   \end{itemize}
\end{itemize}

% ------------------------------------------------------------
\subsection{enumerate}
Dies ist die Standard Nummerierungsliste von \LaTeXe{}. Sie hat einen Abstand zwischen den Einträgen um bei Längeren Zeilen das Lesen zu erleichtern.

\begin{enumerate}
   \item Lorem ipsum dolor sit amet, consectetur adipisicing elit, sed do eiusmod tempor incididunt ut labore et dolore magna aliqua.
%
   \begin{enumerate}
      \item Lorem ipsum dolor sit amet, consectetur adipisicing elit, sed do eiusmod tempor incididunt ut labore et dolore magna aliqua.
      %
      \begin{enumerate}
         \item Lorem ipsum dolor sit amet, consectetur adipisicing elit, sed do eiusmod tempor incididunt ut labore et dolore magna aliqua.
         %
         \item Lorem ipsum dolor sit amet, consectetur adipisicing elit, sed do eiusmod tempor incididunt ut labore et dolore magna aliqua.
         %
         \item Lorem ipsum dolor sit amet, consectetur adipisicing elit, sed do eiusmod tempor incididunt ut labore et dolore magna aliqua.
      \end{enumerate}
      %
      \item Lorem ipsum dolor sit amet, consectetur adipisicing elit, sed do eiusmod tempor incididunt ut labore et dolore magna aliqua.
      %
      \item Lorem ipsum dolor sit amet, consectetur adipisicing elit, sed do eiusmod tempor incididunt ut labore et dolore magna aliqua.
   \end{enumerate}
\end{enumerate}

% ------------------------------------------------------------
%
\subsection{Kompakte Liste}

Die eleganteste Lösung bietet das Paket \package{enumitem} mit der Option `noitemsep'
\begin{itemize}[noitemsep]
\item Diese Umgebung
\item sollte man nur nutzen,
\item wenn die Einträge nicht länger als eine Zeile sind
\end{itemize}

% ------------------------------------------------------------

\subsection{Beliebige Labels}

Dies ist die \env{enumerate} Umgebung des Paketes \package{enumitem} mit der Option [label=\alph{enumi})].

\begin{enumerate}[label=\alph{enumi})]
         \item Lorem ipsum dolor sit amet, consectetur adipisicing elit, sed do eiusmod tempor incididunt ut labore et dolore magna aliqua.
         %
         \item Lorem ipsum dolor sit amet, consectetur adipisicing elit, sed do eiusmod tempor incididunt ut labore et dolore magna aliqua.
\end{enumerate}

% ------------------------------------------------------------
\subsection{description}
%
\begin{description}
\item[Beschreibung] Die Umgebung \env{description} dient zum Erläutern.
\item[Flüsse] Elbe, Weser
\item[Meere] Indischer, Pazifischer Ozean, Mittelmeer
\end{description}

% ============================================================
%\clearpage
\section{Abbildungen}

% ------------------------------------------------------------

\subsection{Normale figure Umgebung}
Nulla malesuada porttitor diam. Donec felis
erat, congue non, volutpat at, tincidunt tristique, libero. Vivamus
viverra fermentum felis. Donec nonummy pellentesque ante. Phasellus
adipiscing semper elit. Proin fermentum massa ac quam. Sed diam
turpis, molestie vitae, placerat a, molestie nec, leo. Maecenas
lacinia.

\begin{figure}[H]
   \centering
   \includegraphics[width=0.3\textwidth]{images/latex}
   \caption[Normale Figure Umgebung]{Dies ist eine lange Abbildungungsbeschriftung. Dies ist eine lange Abbildungungsbeschriftung. Dies ist eine lange Abbildungungsbeschriftung.}
\end{figure}

\subsection{Normale figure Umgebung, zwei Bilder nebeneinander}
\begin{figure}[H]
   \centering
   \begin{minipage}{0.4\linewidth}
   \includegraphics[width=1.0\linewidth]{images/latex}
   \caption[Normale Figure Umgebung]{Dies ist eine lange Abbildungungsbeschriftung. Dies ist eine lange Abbildungungsbeschriftung. Dies ist eine lange Abbildungungsbeschriftung.}
	\end{minipage}\hspace{2em}%\hfill%
   \begin{minipage}{0.4\linewidth}
   \includegraphics[width=1.0\linewidth]{images/latex}
   \caption[Normale Figure Umgebung]{Dies ist eine lange Abbildungungsbeschriftung. Dies ist eine lange Abbildungungsbeschriftung. Dies ist eine lange Abbildungungsbeschriftung.}
	\end{minipage}
\end{figure}

%

% ------------------------------------------------------------
%
\subsection{subcaption in minipages}
	Nam ipsum ligula, eleifend at, accumsan nec, suscipit a,
	ipsum. Morbi blandit ligula feugiat magna. Nunc eleifend consequat
	lorem. Sed lacinia nulla vitae enim. Pellentesque tincidunt purus
	vel magna. Integer non enim. Praesent euismod nunc eu purus. Donec
	bibendum quam in tellus. Nullam cursus pulvinar lectus. Donec et mi.
	Nam vulputate metus eu enim. Vestibulum pellentesque felis eu
	massa.
	
	\begin{figure}[H]
	\begin{minipage}[b]{.5\linewidth}
		\centering
		\includegraphics[width=0.8\linewidth]{images/latex}
		\subcaption{A subfigure}\label{fig:1a}
	\end{minipage}%
	\begin{minipage}[b]{.5\linewidth}
		\centering
		\includegraphics[width=0.8\linewidth]{images/latex}
		\subcaption{Another subfigure}\label{fig:1b}
	\end{minipage}
	\caption{A figure}\label{fig:1}
	\end{figure}

\subsection{subcaption in subfigure Umgebung}
	Nam ipsum ligula, eleifend at, accumsan nec, suscipit a,
	ipsum. Morbi blandit ligula feugiat magna. Nunc eleifend consequat
	lorem. Sed lacinia nulla vitae enim. Pellentesque tincidunt purus
	vel magna. Integer non enim. Praesent euismod nunc eu purus. Donec
	bibendum quam in tellus. Nullam cursus pulvinar lectus. Donec et mi.
	Nam vulputate metus eu enim. Vestibulum pellentesque felis eu
	massa.
	
	\begin{figure}[H]
	\centering
	\begin{subfigure}[b]{.48\linewidth}
		\centering
		\includegraphics[width=0.8\linewidth]{images/latex}
		\caption{A subfigure}\label{fig:2a}
	\end{subfigure}%
	\hfill
	\begin{subfigure}[b]{.48\linewidth}
		\centering
		\includegraphics[width=0.8\linewidth]{images/latex}
		\caption{Another subfigure}\label{fig:2b}
	\end{subfigure}
	\caption{A figure}\label{fig:2}
	\end{figure}

% ------------------------------------------------------------
%
\subsection{captionof}

Pellentesque et lectus a est imperdiet egestas. Fusce tempus facilisis lacus. Morbi porttitor eleifend dolor. Integer ante. Integer ornare. Quisque ac urna. Nam egestas, eros sed tempor tincidunt, dolor ante dapibus tellus, eget dapibus metus augue et lectus.

\begin{center}
   \includegraphics[width=0.3\textwidth]{images/latex}
   \captionof{figure}{Ein Beispiel für ein nichtfließendes Bild mit Caption durch `captionof'}
\end{center}

Nullam sed enim quis enim blandit commodo. Nulla ultricies metus eu lorem. Phasellus tincidunt ullamcorper orci. Cras lectus metus, luctus id, vulputate quis, tempor non, nulla. Nullam ut eros et elit faucibus viverra. Cras eu turpis eu urna ullamcorper feugiat. Aliquam bibendum rhoncus tortor. Etiam eu orci. In congue congue quam. Maecenas nec sem quis mi ultrices luctus. Aliquam et risus eu eros molestie pellentesque. Aenean quis purus sit amet lorem porttitor consectetuer.
%

% ============================================================
\section{Tabellen}

Bei allen hier vorgestellten Tabellen sind Befehle in die Tabellen eingefügt 
die den Stil verändern.

% ------------------------------------------------------------
%
\subsection{Tabelle im \emph{booktabs} Stil}
%
\begin{table}[H]
\tablestyle
\begin{tabular}{lll}
\toprule
   Tabellenkopf &
   Tabellenkopf &
   Tabellenkopf \tabularnewline
\midrule
Inhalt & Inhalt  & Inhalt \tabularnewline
Inhalt & Inhalt  & Inhalt \tabularnewline
Inhalt & Inhalt  & Inhalt \tabularnewline
\bottomrule
\end{tabular}
\end{table}
%


% ------------------------------------------------------------
\subsection{Tabelle ohne Kopf}

\begin{table}[H]
\tablestyle
\begin{tabular}{lll}
\hline
Inhalt & Inhalt  & Inhalt \tabularnewline
Inhalt & Inhalt  & Inhalt \tabularnewline
Inhalt & Inhalt  & Inhalt \tabularnewline
\hline
\end{tabular}
\end{table}

%
\subsection{Tabelle mit tabularx}
Tabelle mit einstellbarer Tabellenbreite
%--------------------------------------------------------
\begin{table}[H]
\tablestyle
\begin{tabularx}{\textwidth}{lXXlX}
\toprule
   Tabellenkopf &
   Tabellenkopf &
   Tabellenkopf &
   Tabellenkopf &
   Tabellenkopf \tabularnewline
\midrule
%
   \textit{Beschreibung} & Inhalt & Inhalt & Inhalt & Inhalt \tabularnewline
   \textit{Beschreibung} & Inhalt & Inhalt & Inhalt & Inhalt \tabularnewline
\bottomrule
\end{tabularx}
\caption{Tabelle mit tabularx}
\end{table}
%

% ------------------------------------------------------------
%
\subsection{Tabelle mit alternatierender Farbe}


\begin{table}[H]
   \tablestyle
   \rowcolors{1}{tablerowcolor}{white!100}
   \begin{tabular}{*{2}{v{0.45\textwidth}}}
   \hline
   \rowcolor{tableheadcolor}
Tabellenkopf &
Tabellenkopf \tabularnewline
\hline
%
 Inhalt  & Inhalt \tabularnewline
 Inhalt  & Inhalt \tabularnewline
 Inhalt  & Inhalt \tabularnewline
 Inhalt  & Inhalt \tabularnewline
 Inhalt  & Inhalt \tabularnewline
 Inhalt  & Inhalt \tabularnewline
\hline
   \end{tabular}
\end{table}
%
% ------------------------------------------------------------


% ============================================================
% \section{Wissenschaftliches}

% ============================================================
\section{Mathematik}

\subsection{amsmath Umgebungen}

Im Folgenden sind einige Beispiel für amsmath Umgebungen aufgelistet. Diese 
sind der Dokumentation `mathmode.pdf' von Herbert Voss entnommen.

\subsubsection{align}
\begin{align}
y      & = d\\
y      & = cx + d\\
y_{12} & = bx^{2} + cx + d\\
y(x)   & = ax^{3} + bx^{2}+ cx + d
\end{align}

\subsubsection{alignat}
\begin{alignat}{3}
abc   &= xxx &&= xxxxxxxxxxxx &&= aaaaaaaaa \\
ab    &= yyyyyyyyyyyyyyy &&= yyyy &&= ab
\end{alignat}

\subsubsection{multline}

\begin{multline}
a+b+c+d+e+f\\
+i+j+k+l+m+n
\end{multline}

\subsubsection{cases}
\begin{equation}
x = \begin{cases}
   0 & \text{if A =...}\\
   1 & \text{if B =...}\\
   x & \parbox{5 cm}{%
\flushleft %
this runs with as much text as you like,
but without an automatic linebreak,
it runs out of page ....} %
\end{cases}
\end{equation}
\newpage
\section{Listings}

Listings können in \LaTeX ~mittels des Paketes \textit{listings} dargestellt werden. Das Paket erlaubt Syntax-Highlighting und unterstützt diverse Programmiersprachen sowie unterschiedliche Skriptsprachen. In Listing \ref{lst:bash} ist beispielhaft Java dargestellt.

\begin{lstlisting}[caption=Java example,label=lst:bash]
/**
* This is a doc comment.
*/
package com.ociweb.jnb.lombok;

import java.util.Date;
import lombok.Data;
import lombok.EqualsAndHashCode;
import lombok.NonNull;

$$@Data
$$@EqualsAndHashCode(exclude={"address","city","state","zip"})
public class Person {
enum Gender { Male, Female }

// another comment

%%@NonNull%% private String firstName;
%%@NonNull%% private String lastName;
%%@NonNull%% private final Gender gender;
%%@NonNull%% private final Date dateOfBirth;

private String ssn;
private String address;
private String city;
private String state;
private String zip;
}
\end{lstlisting}

\section{Literatur}
Latex erlaubt die direkte Einbindung und Zitation von Literatur. Dieses kann entweder über die explizite Definition von \textit{bibitems} erfolgen oder die implizite Definition über \textit{BibTex}-Elemente. In diesem Beispiel findet Alternative \textit{A} Anwendung. Egal welche Technik genutzt wird, muss jedoch der IEEE Zitationsstil eingehalten werden (Siehe ``IEEE Citation Style Guide''). Dieser sieht für jeden denkbaren Veröffentlichungstyp ein bestimmtes Format innerhalb der Literaturliste vor. Innerhalb des Textes werden Zitate dann mittels Indexnummer kenntlich gemacht \cite{zhang1999}.


% ============================================================



%%\begin{document}
%\title{LaTeX Typesetting By Example}
%\author{Phil Farrell\\
%Stanford University School of Earth Sciences}
%\renewcommand{\today}{November 2, 1994}
%\maketitle
\section{LaTeX Typesetting By Example (by Phil Farrell)}
This article demonstrates a basic set of LaTeX formatting commands.
Compare the typeset output side-by-side with the input document.

\subsection{Plain Text}
Type your text in free-format; lines can be as long
or as short
as you wish.
        You can indent         or space out
        your input 
            text in 
                any way you like to highlight the structure
        of your manuscript and make it easier to edit.
LaTeX fills lines and adjusts spacing between words to produce an
aesthetically pleasing result.

Completely blank lines in the input file break your text into
paragraphs.
To change the font for a single character, word, or set of words, 
enclose the word and the font changing command within braces, 
{\em like this}.
A font changing command not enclosed in braces, like the change to \bf 
bold here, keeps that change in effect until the end of the document or
until countermanded by another font switch, like this change back to 
\rm roman. 

\subsection{Displayed Text}
Use the ``quote'' and  ``quotation'' environments for typesetting quoted
material or any other text that should be slightly indented and set off
from the normal text.
\begin{quotation}
The quote and quotation environments are similar, but use different 
settings for paragraph indentation and spacing.

\em When in doubt, consult the manual.
\end{quotation}

So far, I have demonstrated titles, paragraphs, font changes, and
section headings.
Now, I am going to show lists and tables.
\begin{enumerate}
\item 
The ``enumerate'' environment numbers the list elements, like this.

Items in a list can contain multiple paragraphs.
These paragraphs are appropriately spaced and indented according to their
position in the list.
   \begin{itemize}
   \item The ``itemize'' environment sets off list items with ``bullets'',
like this.  Finally, the ``description'' environment lets you put your own
      \begin{description}
      \item[A] label on each item, like this ``A''.
      \item[If the label is long,] the first line of the item text will
be spaced over to the right as needed.
      \end{description}
   \item Of course, lists can be nested, each type up to at least four levels.
One type of list can be nested within another type.
      \begin{itemize}
      \item Nested lists of the same type will change style of numbering 
or ``bullets'' as needed.
      \end{itemize}
   \end{itemize}
\item Don't forget to close off all list environments with the 
appropriate \verb+\end{...}+ command.
Indenting \verb+\begin{...}+, \verb+\item+, and \verb+\end{...}+
commands in the input document according to their nesting level can help 
clarify the structure.
\end{enumerate}

Here is a very simple table showing data lined up in columns.
Notice that I include the table in a ``center'' environment to display
it properly.
The title is created simply as another paragraph in the center environment,
rather than as part of the table itself.

\begin{table}[H]
\centering
\captionabove{Numbers of Computers on Earth Sciences Network, By Type.}
\begin{tabular}{lr}
\hline
Macintosh		& 175\\
DOS/Windows PC	& 60\\
UNIX Workstation or server & 110\\
\hline
\end{tabular}
\end{table}

Here is a more complicated table that has been boxed up, with a multi-column
header and paragraph entries set in one of the columns.
\begin{table}[H]
\centering
\begin{tabular}{lcp{3.5in}}
\toprule
\multicolumn{3}{c}{Places to Go Backpacking}\\ 
\hline
%
Name & Driving Time & Notes\\
     & (hours)      &\\ 
\midrule
Big Basin & 1.5 & 
Very nice overnight to Berry Creek Falls from
either Headquarters or ocean side.\\ 
%
Sunol & 1 & 
Technicolor green in the spring.  Watch out for the cows.\\ 
%
Henry Coe & 1.5 & 
Large wilderness nearby suitable for multi-day treks.\\ 
%
\bottomrule
\end{tabular}
\end{table}

\subsection{Mathematical Equations}
Simple equations, like $x^y$ or $x_n = \sqrt{a + b}$ can be typeset right
in the text line by enclosing them in a pair of single dollar sign symbols.
Don't forget that if you want a real dollar sign in your text, like \$2000,
you have to use the \verb+\$+ command.

A more complicated equation should be typeset in {\em displayed math\/} mode,
like this:
\[
z \left( 1 \ +\  \sqrt{\omega_{i+1} + \zeta -\frac{x+1}{\Theta +1} y + 1} 
\ \right)
\ \ \ =\ \ \  1
\]
The ``equation'' environment displays your equations, and automatically
numbers them consecutively within your document, like this:
\begin{equation}
\left[
{\bf X} + {\rm a} \ \geq\ 
\underline{\hat a} \sum_i^N \lim_{x \rightarrow k} \delta C
\right]
\end{equation}

% in diese Datei gehört der Inhalt des Dokumentes:
\input{content/Auswertung.tex}

\printbibliography
\end{document}
%% Dokument ENDE %%%%%%%%%%%%%%%%%%%%%%%%%%%%%%%%%%%%%%%%%%%%%%%%%%%%%%%%%%

